% =============================================================================
% cap4 §4.1 Visão geral da avaliação (enxuta, sem metalinguagem)
% =============================================================================
\section{Objetivos e postura}
\label{sec:metodologia-visao-geral}

A avaliação visa estimar a viabilidade técnica do pipeline definido no Capítulo~\ref{cap:pipeline}, sem coleta de campo no Campus~2. Em consequência, restringe-se a \textbf{datasets públicos consolidados} de \emph{camera trapping} e re-identificação animal~\cite{beery2018terra,tabak2019nacti,norouzzadeh2018automatically,shinoda2024petface,adam2024wildlifereid10k}: anotações por equipes especializadas, protocolos reproduzíveis e resultados publicados que viabilizam comparação externa. A coleta sistemática no Campus~2, a rotulagem em volume e a construção do catálogo da colônia integram a implementação futura (Capítulo~\ref{cap:conclusao}).

Quatro perguntas orientam o protocolo, uma por fase. \textbf{Fase~II:} o \modeloMD{} discrimina quadros com animal de quadros vazios em material representativo de fauna sinantrópica? \textbf{Fase~III:} o \modeloSN{}, sobre os recortes da Fase~II, discrimina \emph{Felis catus} de classes ecologicamente próximas (procyonidae, didelphidae, felinos silvestres de pequeno porte)? \textbf{Fase~IV:} a combinação de \modeloPetFace{} com o método de partes do corpo de \citeauthor{akbar2025bodypart} (\citeyear{akbar2025bodypart}) produz Rank-$k$ compatível com decisão automática em regime de alta confiança? \textbf{Pipeline:} a execução em cascata mantém a coerência de funil prevista, reduzindo volume e elevando confiança a cada estágio?

A postura é \textbf{operacional}, não comparativa de modelos: os componentes selecionados no Capítulo~\ref{cap:pipeline} operam em modo pré-treinado (\estr{as-is}), com limiares iniciados nos valores padrão dos mantenedores e ajustados, quando necessário, com base nas matrizes de confusão e taxas de abstenção observadas. As métricas adotadas são as consagradas em cada subtarefa (mAP para detecção, $F_1$ e matriz de confusão para classificação, Rank-$k$ e taxa de detecção de novo para re-ID), permitindo situar o sistema em relação à literatura.
